% TEMPLATE-MILC-v3.tex - plantilla maestra UNICA de las guias MILC v3.
%
% La usan las cuatro familias de guias, cada una con su driver:
%   · Grados 6-11          -> scripts/build-guias-g11.py
%   · Bebras               -> scripts/build-guias-bebras.py
%   · Semillero            -> scripts/build-guias-semillero.py
%   · Territorio interior  -> scripts/build-guias-territorio-interior.py
%
% Antes habia tres copias casi identicas (~400 lineas iguales, 6 distintas):
% cada mejora habia que aplicarla tres veces, y en la practica no se hacia
% (el motor de recursos vivia solo en dos de ellas). Lo que varia entre
% familias esta parametrizado en 6 placeholders que cada driver llena
% (se nombran aqui SIN su sintaxis real: el builder reemplaza por texto plano,
% tambien dentro de los comentarios, y un valor multilinea rompe el preambulo):
%
%   HEADER_IZQ        encabezado izquierdo de pagina
%   HEADER_DER        encabezado derecho de pagina
%   PORTADA_ETIQUETA  rotulo sobre el numero grande (GRADO/MOMENTO/MODULO)
%   PORTADA_SERIE     linea de serie de la portada
%   PORTADA_ID        identificador de la guia en la portada
%   PORTADA_CIRCULO   el circulito de grado (°), o vacio si no aplica
%   PDF_TITULO        titulo en texto plano (sin \textbf ni comillas TeX) para
%                     los metadatos del PDF (/Title); lo produce lib_xelatex.texto_pdf
%
% Composicion (auditoria 2026-09-03): las bandas de estacion piden espacio
% (\Needspace*) y prohiben el salto justo despues (\nopagebreak), asi no quedan
% huerfanas al pie; las cajas largas (softbox, drawbox, infoband, puentebox) son
% `breakable`, asi una caja de media pagina ya no arrastra todo a la siguiente
% dejando la anterior al 40 %. Todo texto sobre mostaza o verde va en milcNegro
% (blanco sobre #E5B400 da 1,93:1 y sobre #58A923 2,95:1; ninguno pasa WCAG AA).
% El resto de claves las documenta content/guias/_SCHEMA.md. El script valida
% que no queden placeholders sin reemplazar antes de escribir el archivo final.

% Etiquetado PDF (latex-lab): /Lang, estructura y texto alternativo de las
% figuras, para lectores de pantalla. Debe ir ANTES de \documentclass.
\DocumentMetadata{lang=es-CO,tagging=on}
\documentclass[11pt,letterpaper]{article}
% headheight=24pt: la cabecera derecha lleva dos lineas (identidad de la guia
% y estacion actual). headsep se acorta en la misma medida para que el bloque
% de texto quede exactamente donde estaba con headheight=12pt.
\usepackage[margin=2.0cm,headheight=24pt,headsep=13pt]{geometry}
\usepackage[table]{xcolor}
\usepackage{fontspec}
\usepackage[spanish]{babel}
\usepackage{tikz}
% Librerías TikZ para los diagramas de las guías (bloque `recursos:`):
%  · babel  → babel en español vuelve activos caracteres como «>», y TikZ los usa
%             en sus opciones (>=stealth, ->). Sin esta librería, cualquier
%             diagrama con flechas revienta con «\language@active@arg> has an
%             extra }». Debe ir ANTES de que se dibuje nada.
%  · positioning        → sintaxis «below left=2pt and 3pt».
%  · decorations.pathreplacing → llaves ({brace}) para señalar rangos.
\usetikzlibrary{babel,positioning,decorations.pathreplacing}
\usepackage{graphicx}
% Circuitos: el trabajo de electrónica se hace hoy con micro:bit y sus sensores
% integrados (MakeCode, sin cableado), así que no se necesitan esquemáticos.
% Si algún día entran montajes con protoboard, basta descomentar esta línea:
% los diagramas .tex/.tikz declarados en `recursos:` ya se incrustan solos.
% \usepackage{circuitikz}
\usepackage[most]{tcolorbox}
\usepackage{tabularx}
\usepackage{array}
\usepackage{fancyhdr}
\usepackage{titlesec}
\usepackage[document]{ragged2e}  % bandera a la derecha en todo el cuerpo
\usepackage{enumitem}
\usepackage{microtype}
\usepackage{etoolbox}
\usepackage{needspace}
% hyperref al final del preambulo: metadatos del PDF (titulo, autor, idioma,
% familia), marcadores por estacion (\pdfbookmark) y enlaces sin recuadro.
\usepackage[hidelinks,
  pdftitle={Proyecto integrador comunitario — voz para el cambio social},
  pdfauthor={PhD. Álvaro Cárdenas Orozco},
  pdfsubject={Tecnología e Informática · Grado 8 · Guía 29 · MILC v3},
  pdflang={es-CO},
  bookmarksopen=true,bookmarksnumbered=false]{hyperref}

% Helvetica Neue no trae la flecha U+2192, asi que cada «→» escrito literalmente
% en el YAML se compilaba a NADA, en silencio (462 flechas en 61 guias: rutas del
% tipo «paleta → tipografia → lockup» salian rotas). Mapeamos el caracter a su
% version matematica, que si tiene glifo. Asi el autor puede seguir escribiendo
% «→» comodamente en el YAML.
\usepackage{newunicodechar}
\newunicodechar{→}{\ensuremath{\rightarrow}}
% Mismo problema con los simbolos de marca y vinetas: el «✓» de «una palomita ✓»
% salia como cuadrito vacio en el PDF de todas las guias que lo usaban.
\usepackage{pifont}
\newunicodechar{✓}{\ding{51}}
\newunicodechar{✗}{\ding{55}}
\newunicodechar{✔}{\ding{52}}
\newunicodechar{•}{\textbullet}
\newunicodechar{≥}{\ensuremath{\geq}}
\newunicodechar{≤}{\ensuremath{\leq}}

\setmainfont{Helvetica Neue}
\setsansfont{Helvetica Neue}
\setlength{\parindent}{0pt}
\setlength{\parskip}{6pt}
% Sin particion de palabras: en 6.º-7.º una palabra cortada por guion al final
% de linea («Comuni-cacion») frena la lectura mucho mas que una linea corta.
\hyphenpenalty=10000
\exhyphenpenalty=10000
\pretolerance=2000
\tolerance=3000
\setlength{\emergencystretch}{3em}
\linespread{1.10}
\renewcommand{\arraystretch}{1.25}
\setlist{leftmargin=5mm,itemsep=1mm,topsep=1mm}
\newcolumntype{Y}{>{\RaggedRight\arraybackslash}X}

\definecolor{milcMagenta}{HTML}{E6007E}
\definecolor{milcVino}{HTML}{5A0038}
\definecolor{milcTurquesa}{HTML}{0093A5}
\definecolor{milcVerde}{HTML}{58A923}
\definecolor{milcMostaza}{HTML}{E5B400}
\definecolor{milcOcre}{HTML}{B86600}
\definecolor{milcNegro}{HTML}{241816}
\definecolor{milcGris}{HTML}{F7F7F4}
\definecolor{milcLinea}{HTML}{E8E2DD}
\definecolor{milcGradePrimary}{HTML}{FF6600}
\definecolor{milcGradeDark}{HTML}{9A3E00}
\definecolor{milcGradeSoft}{HTML}{FFE4D0}
\definecolor{milcGradeAccent}{HTML}{CFE7FF}
\definecolor{milcGradeText}{HTML}{FFFFFF}
\color{milcNegro}

\pagestyle{fancy}
\fancyhf{}
% Cabecera de dos lineas: identidad de la guia arriba y, a la derecha y en
% gris, la estacion en curso (\markright desde \milcestacion). Antes de la
% primera estacion la segunda linea va vacia; el \strut mantiene la altura
% para que la linea de arriba no baile entre paginas.
\lhead{\small Tecnología e Informática · Grado 8\\[-1pt]{\footnotesize\strut}}
\rhead{\small Guía 29 · MILC v3\\[-1pt]{\footnotesize\strut\color{milcNegro!65}\rightmark}}
\cfoot{\small\thepage}
\renewcommand{\headrulewidth}{0pt}

% Sobre mostaza (#E5B400) y verde (#58A923) el texto va en milcNegro; sobre el
% resto de la paleta, en blanco. Se usa dentro de fonttitle/fontupper, donde un
% \color posterior manda sobre coltitle.
\newcommand{\milctintasobre}[1]{%
  \ifboolexpr{test{\ifstrequal{#1}{milcMostaza}} or test{\ifstrequal{#1}{milcVerde}}}%
    {\color{milcNegro}}{\color{white}}}

\newtcolorbox{titlebox}[1]{
  enhanced,colback=#1,colframe=#1,arc=7mm,boxrule=0pt,
  left=7mm,right=7mm,top=3mm,bottom=3mm,
  before skip=8pt,after skip=8pt,
  fontupper=\sffamily\bfseries\Large\color{white}
}

\newtcolorbox{softbox}[3]{
  enhanced,breakable,pad at break=3mm,
  colback=#2,colframe=#1,borderline west={4pt}{0pt}{#1},
  arc=2mm,boxrule=.4pt,left=4mm,right=4mm,top=3mm,bottom=3mm,
  before skip=6pt,after skip=8pt,title={#3},coltitle=white,
  colbacktitle=#1,fonttitle=\sffamily\bfseries\milctintasobre{#1},fontupper=\RaggedRight,
  attach boxed title to top left={xshift=3mm,yshift=-2mm},
  boxed title style={arc=2mm, boxrule=0pt}
}

\newtcolorbox{drawbox}[2]{
  enhanced,breakable,pad at break=4mm,
  colback=white,colframe=#1,arc=4mm,boxrule=1pt,
  left=5mm,right=5mm,top=4mm,bottom=4mm,before skip=8pt,after skip=8pt,
  title={#2},coltitle=white,colbacktitle=#1,fonttitle=\sffamily\bfseries\milctintasobre{#1},
  fontupper=\RaggedRight,
  attach boxed title to top left={xshift=4mm,yshift=-2mm},
  boxed title style={arc=3mm, boxrule=0pt}
}

\newtcolorbox{infoband}[1]{
  enhanced,breakable,pad at break=2mm,colback=#1!8,colframe=#1!50,arc=2mm,boxrule=.5pt,
  left=4mm,right=4mm,top=2mm,bottom=2mm,before skip=4pt,after skip=4pt,
  fontupper=\small\RaggedRight
}

% Puente narrativo entre secciones (contrato editorial Conéctate).
% Da continuidad: "Acabas de X. Ahora vas a Y porque Z."
% Visualmente: banda izquierda en color del grado, texto en itálica pequeña.
\newtcolorbox{puentebox}{
  enhanced,breakable,pad at break=2mm,colback=milcGradeSoft,colframe=milcGradeDark,
  borderline west={3pt}{0pt}{milcGradeDark},
  arc=0mm,boxrule=0pt,left=5mm,right=4mm,top=2.5mm,bottom=2.5mm,
  before skip=4pt,after skip=8pt,
  fontupper=\itshape\small\color{milcNegro}\RaggedRight
}
% La flecha va en modo matematico a proposito: Helvetica Neue NO tiene el glifo
% U+2192, asi que \textrightarrow se compilaba a NADA («Missing character: There
% is no → in font Helvetica Neue») y el puente salia sin flecha. En modo
% matematico la flecha viene de la fuente matematica y si se dibuja.
\newcommand{\puente}[1]{%
  \begin{puentebox}%
    \textcolor{milcGradeDark}{$\rightarrow$}\hspace{2mm}#1%
  \end{puentebox}%
}

\newcommand{\linead}[1]{\noindent\color{milcLinea}\rule{#1}{0.5pt}\color{milcNegro}}
\newcommand{\respuesta}[1]{\par\vspace{2mm}\foreach \n in {1,...,#1}{\noindent\color{milcLinea}\rule{\linewidth}{0.5pt}\par\vspace{5mm}}\color{milcNegro}}
\newcommand{\checkbox}{\textcolor{milcGradeDark}{\fbox{\phantom{x}}}\hspace{5pt}}

% ─── Listas aireadas para las guías socioemocionales ────────────────────
% Componer las verificaciones como «\checkbox texto\\» deja el texto que salta
% de línea debajo del cuadrito y apelmaza el bloque. Estos entornos alinean la
% continuación y airean los ítems. Los usa Territorio Interior; cualquier
% familia puede hacerlo.
\newlist{checklista}{itemize}{1}
\setlist[checklista]{
  label=\textcolor{milcGradeDark}{\fbox{\phantom{x}}},
  leftmargin=9mm, labelsep=3.2mm, itemsep=2.4mm,
  topsep=2.5mm, parsep=0pt, partopsep=0pt, before=\RaggedRight
}
\newlist{pasoslista}{enumerate}{1}
\setlist[pasoslista]{
  label=\textcolor{milcGradeDark}{\sffamily\bfseries\arabic*.},
  leftmargin=8mm, labelsep=3mm, itemsep=2.8mm,
  topsep=1mm, parsep=0pt, partopsep=0pt, before=\RaggedRight
}

\newcommand{\phasecard}[4]{
\begin{tcolorbox}[enhanced,colback=#3,colframe=#2,arc=3mm,boxrule=.5pt,height=3.05cm,valign=center,halign=center,left=1.5mm,right=1.5mm,top=2mm,bottom=2mm]
{\sffamily\bfseries\fontsize{22}{24}\selectfont\textcolor{#2}{#1}}\par
{\sffamily\bfseries\small #4}
\end{tcolorbox}
}

% ── Piezas del rediseno 2026 (legibilidad 6.º-11.º) ───────────────────
% Banda de estacion: reemplaza el titlebox macizo. Titulo a la izquierda,
% dato practico (tiempo/modalidad) a la derecha, en una sola linea.
% Nunca huerfana: pide 9 lineas libres (o salta de pagina rellenando con
% \vfill) y prohibe el salto justo despues de la banda. Nueve y no seis:
% la caja partible que sigue necesita ~80pt (titulo adosado, relleno y dos
% lineas) para arrancar en la misma pagina; con menos, tcolorbox la manda
% entera a la siguiente y la banda se queda sola. El penalty 10000 va
% en `after`, ANTES del espacio posterior: un `after skip` (aunque sea 0pt)
% es un \vskip tras la caja, y ese glue es un punto de corte legal que un
% \nopagebreak escrito despues de \end{tcolorbox} ya no cierra. Cada banda
% deja marcador PDF y actualiza la cabecera derecha.
\newcounter{milcestacion}
\newcommand{\milcestacion}[2]{%
\Needspace*{9\baselineskip}%
\stepcounter{milcestacion}%
\pdfbookmark[1]{#2}{milcestacion\themilcestacion}%
\markright{#2}%
\begin{tcolorbox}[enhanced,colback=#1,colframe=#1,arc=2mm,boxrule=0pt,
  left=5mm,right=5mm,top=2.6mm,bottom=2.6mm,before skip=11pt,
  after={\par\nopagebreak\vskip7pt}]
{\sffamily\bfseries\fontsize{15}{18}\selectfont\milctintasobre{#1}#2}
\end{tcolorbox}}

% Tarjeta de paso numerada: sustituye las tablas «Pilar / Aplicacion», que
% eran formato de consulta docente, no de estudiante. El numero va en un
% circulo sobre el borde para que la secuencia se lea de un vistazo.
\newcommand{\milcpaso}[3]{%
\Needspace{4\baselineskip}%
\begin{tcolorbox}[enhanced,colback=white,colframe=milcLinea,arc=1.5mm,boxrule=.9pt,
  left=14mm,right=4mm,top=3mm,bottom=3mm,before skip=4pt,after skip=4pt,
  fontupper=\RaggedRight,
  overlay={\node[anchor=north west,circle,fill=#2,text=white,inner sep=0pt,
    minimum size=7.4mm,font=\sffamily\bfseries\small]
    at ([xshift=3.6mm,yshift=-3.2mm]frame.north west) {#1};}]
#3
\end{tcolorbox}}

% Casilla marcable de tamano util con lapiz (la anterior era \fbox{\phantom{x}},
% de alto variable segun el texto que la seguia).
\newcommand{\milcmarca}{\framebox[3.8mm]{\rule{0pt}{3.4mm}}}

% Fila marcable de lista suelta (checks de la estacion 1).
\newcommand{\milccheckrow}[1]{%
\par\vspace{1.8mm}\noindent
\makebox[6.4mm][l]{\raisebox{-.55mm}{\textcolor{milcNegro}{\milcmarca}}}%
\parbox[t]{\dimexpr\linewidth-6.4mm\relax}{\RaggedRight #1}\par}

% Celda marcable dentro de la rubrica (tabla de autoevaluacion). Misma
% sangria francesa, pero pensada para vivir dentro de una columna X.
\newcommand{\milcopcion}[1]{%
\makebox[5.2mm][l]{\raisebox{-.55mm}{\textcolor{milcNegro}{\milcmarca}}}%
\parbox[t]{\dimexpr\linewidth-5.2mm\relax}{\RaggedRight #1}}

% ── Recursos visuales de la guía ──────────────────────────────────────
% Declarados en el bloque `recursos:` del YAML y emitidos por el builder.
% \guiaFigura   → imagen rasterizada o PDF (foto, carta celeste, montaje).
% \guiaDiagrama → fuente TikZ/circuitikz (.tex) incrustada como vector:
%                 es la vía para circuitos y esquemáticos editables.
% [#1] = texto alternativo (alt) · #2 = ruta relativa al .tex · #3 = pie
% El alt llega al PDF etiquetado (/Alt) y lo lee el lector de pantalla.
\newcommand{\guiaFigura}[3][]{%
\begin{center}
\includegraphics[width=0.88\linewidth,height=0.38\textheight,keepaspectratio,alt={#1}]{#2}\\[1.5mm]
{\footnotesize\itshape\textcolor{milcNegro!75}{#3}}
\end{center}
}

\newcommand{\guiaDiagrama}[2]{%
\begin{center}
\input{#1}\\[1.5mm]
{\footnotesize\itshape\textcolor{milcNegro!75}{#2}}
\end{center}
}

\begin{document}
\thispagestyle{empty}

% ── Portada ───────────────────────────────────────────────────────────
% Antes: un rectangulo de color a pagina completa. Se veia bien en pantalla,
% pero en la fotocopiadora del colegio se comia el toner y salia gris sucio.
% Ahora la banda ocupa el tercio superior y el resto va en blanco.
\begin{tikzpicture}[remember picture,overlay]
  \path[fill=milcGradePrimary]
    (current page.north west) rectangle ([yshift=-10.6cm]current page.north east);

  \node[anchor=north west,text=milcGradeText] at ([xshift=2.0cm,yshift=-1.55cm]current page.north west)
    {\sffamily\bfseries\fontsize{12}{14}\selectfont GRADO};
  \node[anchor=north west,text=milcGradeText,opacity=.85] at ([xshift=2.0cm,yshift=-2.35cm]current page.north west)
    {\sffamily\bfseries\fontsize{8.5}{10}\selectfont SERIE GUÍAS · TECNOLOGÍA E INFORMÁTICA};
  \node[anchor=north east,text=milcGradeText,opacity=.85,align=right] at ([xshift=-2.0cm,yshift=-1.55cm]current page.north east)
    {\sffamily\bfseries\fontsize{9}{12}\selectfont I.E. Sor María Juliana\\Cartago, Valle del Cauca};

  \node[anchor=west,text=milcGradeText] (gradonum) at ([xshift=1.85cm,yshift=-7.1cm]current page.north west)
    {\sffamily\bfseries\fontsize{112}{112}\selectfont 8};
  % El circulito de grado (°) solo aplica a las guias de grado («11°»); en
  % Bebras y Semillero el numero grande es un momento/modulo, no un grado.
  % Se posiciona relativo al nodo `gradonum`, no en coordenadas absolutas.
  \draw[milcGradeText,line width=1.6pt]
    ([xshift=2.0mm,yshift=-4.5mm]gradonum.north east) circle (.15cm);

  \node[anchor=west,text=milcGradeText,text width=11.4cm,align=left] at ([xshift=1.5cm,yshift=0.5cm]gradonum.east)
    {
     {\sffamily\bfseries\fontsize{9}{11}\selectfont\MakeUppercase{Período 3 · Multimedia, narrativa y ciberseguridad}}\\[3mm]
     {\sffamily\bfseries\fontsize{21}{25}\selectfont\setlength{\baselineskip}{27pt}Proyecto integrador\\[4pt]comunitario ---\\[4pt]voz para el cambio}\\[3.5mm]
     {\sffamily\bfseries\fontsize{15}{18}\selectfont Guía 29-8-TIC}
    };
\end{tikzpicture}

\vspace*{9.4cm}
\vfill

{\sffamily\bfseries\fontsize{8}{10}\selectfont\textcolor{milcGradeDark}{LO QUE TE LLEVAS HOY}}

\begin{tcolorbox}[enhanced,colback=white,colframe=milcNegro,arc=2mm,boxrule=1.6pt,
  left=5mm,right=5mm,top=4mm,bottom=4mm,before skip=3pt,after skip=12pt,
  fontupper=\RaggedRight\fontsize{11}{15.5}\selectfont]
Un proyecto multimedia sobre un problema real de tu barrio, tu cuadra o tu colegio, con cuatro piezas: la pieza principal, en video corto, audio o infografía, con la estructura de cinco partes; la entrevista a la persona afectada, con su consentimiento documentado; una propuesta de tres acciones con responsables nombrados, no «el gobierno»; y la bitácora del proceso, que registra a quién contactaste, qué dejaste fuera y si le mostraste el resultado antes de publicarlo.
\end{tcolorbox}

\begin{tcolorbox}[enhanced,colback=milcGris,colframe=milcGris,arc=2mm,boxrule=0pt,
  left=5mm,right=5mm,top=3.5mm,bottom=3.5mm,after skip=12pt]
{\sffamily\bfseries\fontsize{8}{10}\selectfont\textcolor{milcNegro!55}{ESTA GUÍA ES DE}}\\[3mm]
\begin{tabularx}{\linewidth}{X p{3.2cm} p{3.2cm}}
\linead{\linewidth} & \linead{3.0cm} & \linead{3.0cm}\\[-1mm]
{\fontsize{8.5}{10}\selectfont\textcolor{milcNegro!55}{Nombre completo}} &
{\fontsize{8.5}{10}\selectfont\textcolor{milcNegro!55}{Grupo}} &
{\fontsize{8.5}{10}\selectfont\textcolor{milcNegro!55}{Fecha}}\\
\end{tabularx}
\end{tcolorbox}

\vfill

\noindent\textcolor{milcLinea}{\rule{\linewidth}{1pt}}\\[2mm]
{\sffamily\bfseries\fontsize{9.5}{12}\selectfont Autor: Dr. Álvaro Cárdenas Orozco}\\[1mm]
{\sffamily\fontsize{8.5}{11}\selectfont\textcolor{milcNegro!55}{Metodología MILC · apertura ancestral · pensamiento computacional · triángulo Dussel–estoicismo–Floridi}}

\newpage

\section*{Apertura: Proyecto integrador comunitario --- voz para el cambio social}

\begin{softbox}{milcGradeDark}{milcGradeSoft}{Saber ancestral}
Desde 1958, Colombia tiene una institución donde los vecinos deciden juntos: la \textbf{Junta de Acción Comunal} (Ley 19 de 1958). Hoy la rige la Ley 2166 de 2021, y ahí están las reglas exactas del juego. En una vereda bastan veinte afiliados para constituirla. La asamblea se reúne por lo menos tres veces al año. No se puede ni siquiera deliberar con menos del veinte por ciento de los miembros. Para cambiar los estatutos se necesitan dos terceras partes. Y hay un dato que casi nadie de tu edad conoce: se puede ser afiliado desde los catorce años. Eso significa que en tu barrio, en Cartago o donde vivas, ya existe una mesa donde tu propuesta puede entrar. La \textbf{cara de exclusión} está documentada: la asamblea es también donde el clientelismo se vuelve concreto. Quien tiene la relación personal con el alcalde consigue la obra (Líppez-De Castro et al., 2020). Deliberar bien no garantiza incidir. Hoy vas a hacer una propuesta que nombre a quién le corresponde responder.\par\smallskip{\footnotesize\textcolor{milcNegro!70}{\textbf{Fuente.} Congreso de la República de Colombia. (2021, 18 de diciembre). \emph{Ley 2166 de 2021. Por la cual se establece el régimen de los organismos de acción comunal\ldots} Diario Oficial.}}
\end{softbox}

\begin{softbox}{milcTurquesa}{milcGris}{Saber contemporáneo}
Un proyecto comunitario tiene cuatro propiedades que no se negocian. \textbf{Problema real}: observado por ti, no copiado de internet. \textbf{Voz de quien lo sufre}: con nombre, oficio, contexto y consentimiento explícito. \textbf{Propuesta concreta}: tres acciones con responsables nombrables, como la JAC del barrio o la coordinación del colegio. \textbf{Pieza producida}: terminada y lista para circular, no planeada. La pieza ética central es el \textbf{consentimiento}. Antes de entrevistar a alguien le explicas cuatro cosas. Que es un trabajo escolar. Que su testimonio irá en una pieza que se mostrará en clase y quizá en redes. Que puede negarse o pedir cambios. Que puede pedir que se retire su nombre después. Y queda registrado: un mensaje, un audio, una firma. Sin ese registro, el proyecto puede hacer daño real a la persona que quiso ayudarte.
\end{softbox}

\begin{softbox}{milcMostaza}{milcGris}{Pregunta-puente}
\textit{La ley dice que se puede ser afiliado a la junta desde los catorce años, y casi nadie lo usa. Tu propuesta va a nombrar a alguien que responda. ¿A quién, y por qué a esa persona?}
\end{softbox}

\begin{softbox}{milcVerde}{milcGris}{Saber y saber hacer}
Al terminar podrás: (1) \textbf{identificar} problemas reales de tu entorno con una persona afectada a la que puedas contactar esta semana; (2) \textbf{analizar} las cuatro propiedades del proyecto comunitario y preparar el consentimiento; (3) \textbf{crear} la pieza, la propuesta de tres acciones con responsables y la bitácora del proceso.
\end{softbox}



\small
\begin{tabularx}{\textwidth}{>{\bfseries}p{3.0cm}Y}
\rowcolor{milcGradePrimary!12}Autor & Dr. Álvaro Cárdenas Orozco\\
DBA & Producción multimedia ética y comportamiento responsable en entornos digitales (MEN, T\&I)\\
Referentes & Dussel (estética de la liberación) · Ley 1336 · Floridi (infoesfera)\\
Producto final & Un proyecto multimedia sobre un problema real de tu barrio, tu cuadra o tu colegio, con cuatro piezas: la pieza principal, en video corto, audio o infografía, con la estructura de cinco partes; la entrevista a la persona afectada, con su consentimiento documentado; una propuesta de tres acciones con responsables nombrados, no «el gobierno»; y la bitácora del proceso, que registra a quién contactaste, qué dejaste fuera y si le mostraste el resultado antes de publicarlo.\\
\end{tabularx}
\normalsize

\puente{Todo el año apunta aquí. Datos en el primer periodo, lógica y sensores en el segundo, multimedia y ética en el tercero. Hoy se juntan sobre un problema que tú viste pasar. En la sesión 10 lo sustentas en público y cierras el grado.}

\milcestacion{milcGradeDark}{Ruta MILC: así vamos a aprender}
\begin{tabularx}{\textwidth}{>{\centering\arraybackslash}X>{\centering\arraybackslash}X>{\centering\arraybackslash}X>{\centering\arraybackslash}X}
\phasecard{1}{milcVerde}{milcGris}{Escucha\\[-1mm]\small Observo, escucho y pregunto.} &
\phasecard{2}{milcTurquesa}{milcGris}{Entender\\[-1mm]\small Sistematizo y ordeno.} &
\phasecard{3}{milcMagenta}{milcGris}{Hacer\\[-1mm]\small Construyo y aplico.} &
\phasecard{4}{milcVino}{milcGris}{Cierre\\[-1mm]\small Evalúo y reflexiono.}
\end{tabularx}


\puente{Antes de la teoría vas a hacer un inventario de tu entorno. Cinco problemas que hayas visto, no leído. Con una persona afectada en cada uno, y la pregunta clave: ¿puedo hablar con ella esta semana?}

\milcestacion{milcVerde}{Estación 1 de 3 · Escucha}

\begin{softbox}{milcVerde}{milcGris}{Escena para escuchar}
\textbf{Actividad 1 · IDENTIFICA --- Inventario del entorno} (15 min · individual). (1) Anota cinco problemas de tu barrio, tu cuadra o tu colegio que hayas visto pasar, no leído en noticias. (2) Para cada uno, escribe a quién afecta y si conoces a esa persona. (3) Marca en cuáles podrías hablar con la persona afectada esta semana. (4) Elige uno como finalista y escribe su nombre completo. Ejemplos válidos: un paso peatonal sin semáforo, un aula con goteras, la basura de una esquina concreta. También un negocio del barrio que pierde clientes por no tener catálogo digital.
\end{softbox}

{\sffamily\bfseries\fontsize{8}{10}\selectfont\textcolor{milcNegro!55}{MARCA CUANDO LO HAYAS HECHO}}
\milccheckrow{Anoté cinco problemas que vi pasar, no leídos.}
\milccheckrow{Escribí a quién afecta cada uno y si lo conozco.}
\milccheckrow{Elegí uno con una persona a la que puedo contactar esta semana.}

\begin{infoband}{milcVerde}


\textbf{Lo que va al cuaderno (Actividad 1 · IDENTIFICA).} \textbf{Título:} «Actividad 1 --- Inventario del entorno». \textbf{Formato:} tabla de 5 filas y 3 columnas (problema visto / a quién afecta / puedo contactarla esta semana), más el finalista con nombre. \textbf{Extensión:} un tercio de página. \textbf{Sabes que terminaste cuando} tu finalista tiene una persona real detrás y una forma de contactarla.
\end{infoband}

\puente{Ya tienes candidatos. Ahora vas a ordenar las cuatro propiedades que hacen que esto sea un proyecto y no un ejercicio, y a preparar cómo vas a pedir el consentimiento.}

\milcestacion{milcTurquesa}{Estación 2 de 3 · Entender}

\begin{softbox}{milcTurquesa}{milcGris}{Lo que vas a entender}
\textbf{Actividad 2 · ANALIZA --- Las 4 propiedades y el consentimiento} (30 min · parejas). (1) Con tu pareja, revisen los dos finalistas contra las cuatro propiedades y anoten cuál cumple mejor. (2) Escriban las tres acciones de la propuesta, cada una con el nombre del responsable. (3) Redacten el texto exacto con el que van a pedir el consentimiento, con las cuatro cosas que hay que explicar. (4) Escriban las cinco preguntas de la entrevista y marquen cuál es la que de verdad quieren que les contesten.
\end{softbox}

\milcpaso{1}{milcTurquesa}{\textbf{Las cuatro propiedades.} \textbf{Problema real}: lo viste tú. \textbf{Voz de quien lo sufre}: con nombre, oficio y contexto. \textbf{Propuesta concreta}: tres acciones, cada una con quién debería hacerla. \textbf{Pieza producida}: terminada, no planeada.}
\milcpaso{2}{milcTurquesa}{\textbf{El responsable tiene nombre.} «El gobierno debería» no es una propuesta: es una queja. «La JAC del barrio», «la coordinación del colegio», «los vecinos de la cuadra impar» sí lo son. Si no sabes a quién le corresponde, esa averiguación es parte del trabajo.}
\milcpaso{3}{milcTurquesa}{\textbf{El consentimiento, en cuatro frases.} Es un trabajo del colegio. Lo que usted diga irá en una pieza que se mostrará en clase y quizá en redes. Puede negarse o pedir cambios. Puede pedir después que se quite su nombre. Y queda registrado en un mensaje, un audio o una firma.}
\milcpaso{4}{milcTurquesa}{\textbf{La bitácora es la firma del oficio.} Anota cuándo y cómo contactaste, cuándo fue la entrevista, qué decisiones de edición tomaste, qué dejaste fuera y si le mostraste el resultado antes de publicarlo. Sin ese registro no se puede saber si el trabajo se hizo bien.}

\begin{drawbox}{milcTurquesa}{El proyecto comunitario, en una mirada}
\textbf{4 propiedades:} problema real / voz con nombre / propuesta con responsables / pieza producida. \\[1mm] \textbf{Consentimiento:} qué es, dónde saldrá, puede negarse, puede retirarse. \\[1mm] \textbf{Propuesta:} 3 acciones, 3 responsables. \\[1mm] \textbf{Bitácora:} contacto, entrevista, edición, qué quedó fuera.
\end{drawbox}

\begin{softbox}{milcMostaza}{milcGris}{Errores comunes y su corrección}
(1) \textbf{Elegir un problema enorme y lejano}: el cambio climático no cabe; la basura de tu esquina sí. (2) \textbf{Denunciar sin proponer}: señalar el problema y no decir qué hacer con él. (3) \textbf{Poner de responsable a «el gobierno»}: sin nombre, nadie responde. (4) \textbf{Entrevistar sin consentimiento}: usar el testimonio de alguien que no supo dónde iba a salir. (5) \textbf{Producir solo para la nota}: si la pieza no puede circular, el proyecto se quedó en ejercicio.
\end{softbox}

\begin{infoband}{milcTurquesa}


\textbf{Lo que va al cuaderno (Actividad 2 · ANALIZA).} \textbf{Título:} «Actividad 2 --- Las 4 propiedades y el consentimiento». \textbf{Formato:} las cuatro propiedades aplicadas al finalista, las tres acciones con su responsable, el texto exacto del consentimiento y las cinco preguntas. \textbf{Extensión:} media página. \textbf{Sabes que terminaste cuando} cada acción tiene un responsable con nombre.
\end{infoband}

\puente{Tienes el plan. Toca ejecutar: contactar, entrevistar, producir, proponer y dejar constancia de todo en la bitácora.}

\milcestacion{milcMagenta}{Estación 3 de 3 · Hacer}

\begin{softbox}{milcMagenta}{milcGris}{Cómo vas a construir tu producto}
\textbf{Actividad 3 · CREA --- El proyecto ejecutado} (25 min · parejas). (1) Contacta a la persona y pide el consentimiento con el texto que escribieron. (2) Haz la entrevista en audio o video, con las cinco preguntas. (3) Produce la pieza con la estructura de cinco partes. (4) Redacta la propuesta de tres acciones con sus responsables. (5) Llena la bitácora y, si puedes, muéstrale la pieza a la persona antes de publicarla. La sesión alcanza para arrancar; la producción se termina durante la semana.
\end{softbox}

\milcpaso{1}{milcMagenta}{\textbf{Tu entrega tiene cuatro piezas.} La pieza multimedia. La entrevista con el consentimiento registrado. La propuesta de tres acciones con responsables. La bitácora del proceso.}
\milcpaso{2}{milcMagenta}{\textbf{La estructura de la pieza, en cinco partes.} Qué problema es y a quién afecta, en veinte segundos. La voz de la persona, con una frase textual, en un minuto. El dato o el contexto que lo dimensiona, en media. Las tres acciones con sus responsables, en otra media. Y el cierre, con lo que le pides a quien mira.}
\milcpaso{3}{milcMagenta}{\textbf{Las cinco preguntas de la entrevista.} Qué le pasa. Desde cuándo. Qué ha intentado. Qué le ayudaría. Qué cambiaría si tuviera apoyo. La cuarta es la que casi nadie hace, y es la que da la propuesta.}
\milcpaso{4}{milcMagenta}{\textbf{Devolverle la pieza a quien te habló.} Es lo que separa este trabajo de una tarea. Si la persona ve el resultado antes de que circule, puede corregir lo que quedó mal dicho, que es exactamente el derecho que le prometiste.}

\begin{drawbox}{milcMagenta}{Antes de entregar}
\checkbox Problema visto por ti, con una persona afectada contactada. \\[1mm] \checkbox Consentimiento registrado en mensaje, audio o firma. \\[1mm] \checkbox Entrevista hecha con las cinco preguntas. \\[1mm] \checkbox Pieza con la estructura de cinco partes. \\[1mm] \checkbox Tres acciones, cada una con su responsable nombrado. \\[1mm] \checkbox Bitácora con lo que quedó fuera. \\[1mm] \checkbox Copia entregada a la persona entrevistada.
\end{drawbox}

\begin{softbox}{milcMagenta}{milcGris}{Plantilla de trabajo}
\textbf{Problema.} \textit{Visto por mí, en tal lugar.} \\[1mm] \textbf{Persona.} \textit{Nombre, oficio, consentimiento registrado.} \\[1mm] \textbf{Pieza.} \textit{Formato y cinco partes.} \\[1mm] \textbf{Propuesta.} \textit{Tres acciones, tres responsables.} \\[1mm] \textbf{Bitácora.} \textit{Contacto, edición, qué quedó fuera.}
\end{softbox}

\begin{infoband}{milcMagenta}


\textbf{Lo que va al cuaderno (Actividad 3 · CREA).} \textbf{Título:} «Actividad 3 --- El proyecto ejecutado». \textbf{Formato:} la referencia a la pieza, el registro del consentimiento, la propuesta con responsables y la bitácora. \textbf{Extensión:} media página. \textbf{Sabes que terminaste cuando} le entregaste una copia a la persona que entrevistaste.
\end{infoband}

\puente{Lo que entregas hoy: la pieza, la entrevista con consentimiento, la propuesta y la bitácora. La prueba de calidad: la persona entrevistada podría ver la pieza y reconocer lo que dijo.}

\milcestacion{milcMostaza}{Producto de la sesión}
\textbf{Pieza multimedia + entrevista con consentimiento + propuesta con responsables + bitácora}.
\begin{softbox}{milcMostaza}{milcGris}{Criterios de calidad}
(1) El problema lo viste tú y tiene una persona afectada con nombre. (2) El consentimiento quedó registrado en mensaje, audio o firma. (3) La pieza tiene las cinco partes y está terminada, no planeada. (4) Cada una de las tres acciones nombra a quién le corresponde hacerla. (5) La bitácora dice qué quedó fuera y si le mostraste la pieza a la persona.
\end{softbox}

\puente{Antes de cerrar, mira lo que hiciste desde cinco lados. Una denuncia sin propuesta deja el problema donde estaba y de paso expone a quien te habló.}

\milcestacion{milcVino}{Evaluación liberadora · cinco dimensiones}

\begin{softbox}{milcVino}{milcGris}{Sentido de la evaluación}
Evaluar no es solo revisar una nota. Es reconocer qué aprendí, qué cambió en mí, qué emociones manejé, cómo me relaciono con la ciudad y la comunidad, y qué le dejaré a quien viene detrás.
\end{softbox}

\textbf{Mi reflexión liberadora en cinco dimensiones.} Completa cada frase iniciada:

\small
\begin{tabularx}{\textwidth}{>{\bfseries\RaggedRight\arraybackslash}p{3.4cm}Y>{\RaggedRight\arraybackslash}p{4.6cm}}
\rowcolor{milcVino}\color{white}Dimensión & \color{white}\bfseries Mi respuesta (frame starter) & \color{white}\bfseries Algo que descubrí\\
1. Personal & Lo que más cambió en mí fue \linead{6.5cm} & \linead{4.2cm}\\
2. Emocional & La emoción más fuerte fue \linead{2.5cm} y la regulé \linead{3cm} & \linead{4.2cm}\\
3. Ciudadana & En democracia esto importa porque \linead{6cm} & \linead{4.2cm}\\
4. Local (Cartago) & En mi barrio sería útil para \linead{6.5cm} & \linead{4.2cm}\\
5. Intergeneracional & A mi abuela/abuelo le diría \linead{6.5cm} & \linead{4.2cm}\\
\end{tabularx}
\normalsize

\begin{infoband}{milcVino}
\textbf{Para la próxima sustentación.} Vuelve a esta tabla en seis meses. Verás que las respuestas cambian — no porque lo aprendido cambie, sino porque tú cambias. La evaluación liberadora es una conversación contigo a través del tiempo.
\end{infoband}


\puente{Para terminar, tres ideas sobre escuchar a alguien y hacer algo con eso, contadas con palabras de esta guía: Dussel sobre el otro que reclama, Marco Aurelio sobre entrar en lo que dice quien te habla, Floridi sobre saber preguntar.}

\milcestacion{milcGradeDark}{Tres ideas para cerrar · Dussel, un estoico y Floridi}

\begin{softbox}{milcVino}{milcGris}{Enrique Dussel · lente del nosotros}
\textit{El otro se revela de verdad cuando aparece desde fuera del sistema y reclama algo que le corresponde.}

{\footnotesize\itshape\color{milcNegro!70}--- idea tomada de Enrique Dussel, contada con palabras de esta guía. No es una cita textual.}

La persona que entrevistas no es el tema de tu pieza: es quien reclama. Por eso la propuesta va después de su voz y no antes, y por eso nombra a alguien que debería responderle.

\textbf{¿Mi pieza deja que esa persona reclame, o solo la usa para ilustrar lo que yo quería decir?}
\end{softbox}
\linead{\linewidth}\\[1mm]
\linead{\linewidth}

\begin{softbox}{milcMostaza}{milcGris}{Estoicismo · Marco Aurelio · cuidado interior}
\textit{Acostúmbrate a atender lo que dice el otro y, hasta donde puedas, métete en lo que está pensando.}

{\footnotesize\itshape\color{milcNegro!70}--- idea tomada de Marco Aurelio, contada con palabras de esta guía. No es una cita textual.}

En una entrevista se nota enseguida quién escucha y quién espera su turno para la siguiente pregunta. La cuarta pregunta, qué le ayudaría, solo funciona si de verdad escuchaste las tres anteriores.

\textbf{¿Pregunté lo que llevaba escrito, o pregunté lo que apareció mientras me hablaba?}
\end{softbox}
\linead{\linewidth}\\[1mm]
\linead{\linewidth}

\begin{softbox}{milcTurquesa}{milcGris}{Luciano Floridi · lente de la infoesfera}
\textit{Gana quien sabe preguntar y responder, y por eso sabe qué datos vale la pena recoger y cuidar.}

{\footnotesize\itshape\color{milcNegro!70}--- idea tomada de Luciano Floridi, contada con palabras de esta guía. No es una cita textual.}

Tu proyecto no necesita muchos datos: necesita los que sostienen la propuesta. Una fecha, una frecuencia, un número de personas afectadas. Lo demás sobra y distrae de lo que pediste.

\textbf{¿Qué dato de mi pieza sostiene de verdad la propuesta, y cuál está solo de adorno?}
\end{softbox}
\linead{\linewidth}\\[1mm]
\linead{\linewidth}

\begin{infoband}{milcNegro}
\textbf{Nota para el docente.} Las tres ideas están contadas con palabras de esta guía a partir del pensamiento de cada autor; no son citas textuales. De 9.º en adelante se trabajan con la obra y su referencia.
\end{infoband}

\milcestacion{milcVino}{Mi autoevaluación · marca una casilla por fila}

{\small
\setlength{\extrarowheight}{2.2pt}
\begin{tabularx}{\textwidth}{>{\RaggedRight\arraybackslash\bfseries}p{3.0cm}YYY}
\rowcolor{milcVino}\color{white}\bfseries Criterio & \color{white}\bfseries Lo logré & \color{white}\bfseries En proceso & \color{white}\bfseries Necesito apoyo\\[.6mm]
Comprensión del tema & \milcopcion{Explico las ideas con mis palabras.} & \milcopcion{Reconozco las ideas, falta precisión.} & \milcopcion{Necesito repasar lo básico.}\\[.8mm]
\rowcolor{milcGris}Pensamiento computacional & \milcopcion{Usé los pilares al armar mi producto.} & \milcopcion{Usé algunos pilares.} & \milcopcion{Necesito releer la estación 2.}\\[.8mm]
Aplicación y producto & \milcopcion{Mi producto cumple los criterios.} & \milcopcion{Está armado, con ajustes pendientes.} & \milcopcion{Aún está en construcción.}\\[.8mm]
\rowcolor{milcGris}Reflexión 5 dimensiones & \milcopcion{Respondí cada una con honestidad.} & \milcopcion{Respondí algunas con detalle.} & \milcopcion{Mis respuestas son muy generales.}\\[.8mm]
Triángulo de pensamiento & \milcopcion{Conecté las 3 voces con mi trabajo.} & \milcopcion{Conecté al menos una.} & \milcopcion{No las relaciono con mi proceso.}\\[.8mm]
\end{tabularx}}

\vspace{3mm}
\textbf{Hoy entendí que:}\\[1mm]
\linead{\linewidth}\\[1mm]
\linead{\linewidth}

\textbf{Mi compromiso para la próxima sesión es:}\\[1mm]
\linead{\linewidth}\\[1mm]
\linead{\linewidth}

\begin{softbox}{milcMostaza}{milcGris}{Cita de despedida · Marco Aurelio}
\textit{``El obstáculo en el camino se vuelve el camino. Lo que estorba al camino se vuelve camino.''}\\[1mm]
Cada error de hoy, bien observado, se convierte en pieza del camino que pisarás mañana.
\end{softbox}

\small
\begin{tabularx}{\textwidth}{>{\bfseries}p{3.6cm}Y}
\rowcolor{milcGradeDark!12}Nombre completo: & \linead{10cm}\\
Fecha: & \linead{4cm} \quad Curso/grupo: \linead{4cm}\\
Mi firma: & \linead{9cm}\\
\end{tabularx}
\normalsize

\begin{infoband}{milcGradeDark}
\textbf{Para la próxima sesión.} Llega con la pieza, el consentimiento, la propuesta y la bitácora. Si puedes, hazle llegar una copia a la persona que entrevistaste. En la sesión 10 sustentas este proyecto en público y armas el portafolio del año.
\end{infoband}

\end{document}
